\section{Intellectuality and Organic Logic}

\begin{quotex}
As a negative impulse, conservatism is based on a certain distrust of human nature, believing that the immediate impulses of the heart and visions of the brain are likely to be misleading guides. But with this distrust of human nature is closely connected another and more positive factor of conservatism --- its trust in the controlling power of the imagination.
\flright{\textsc{Paul Elmer More}}

The destructive and monstrous opinion that no one, or few, should philosophize, has much invaded the minds of almost everybody. As if it were absolutely nothing to have \textbf{the causes of things},
\textbf{the ways of nature}, \textbf{the reason of the universe}, \textbf{the counsel of God}, \textbf{the mysteries of heaven and earth} very certain before our eyes and hands, unless someone could derive some benefit from it or acquire profit for himself.
\flright{\textsc{Pico della Mirandola}, \emph{Oration on the Dignity of Man}}
\end{quotex}


In \textbf{Valentin Tomberg}'s meditation on the 18th Arcanum of the Tarot, \emph{The Moon} is identified with ``materialistic intellectuality''. This then eclipses the ``creative light of consciousness'' for human intelligence. This materialistic intellectuality seeks the causes of things in the ``least developed and most primitive'' aspects of life. It has regard for the life of the mind only insofar as it brings material benefit. That is why we are far less sanguine than Pico about the possibility that more than the few will ever philosophize in Pico's sense.

A step above that, and probably more sinister, is to follow every ``impulse of the heart'' or any idea that pops into one's mind as guides to life. This ``magical enchantment'' to the lower forces prevents the moral imagination from cognizing higher realities. This sequence that begins with sense life, which then excites the passions, and finally directing the mind, is a deformation of the normal sequence: viz., thought awakens feeling and directs the will. As the Cabala describes, the emanation from the One to the manifest world descends through the three soul activities of thought, feeling and will in two pillars. (The return path is described in The Middle Pillar.) The left pillar refers to natural knowledge and the right pillar to wisdom, or revealed knowledge as summarized in Table \ref{tab:pillars}:

\begin{table}[h]\centering\small
\begin{tabular}{p{.3\textwidth}p{.3\textwidth}p{.3\textwidth}}\toprule
\textbf{Soul Activity} & \textbf{Left or Feminine Pillar (Severity)} & \textbf{Right or Masculine Pillar (Mercy)} \\\midrule
Thinking & Reason or pure intellect & Wisdom or Revelation \\
Feeling & Justice & Magnanimity \\
Willing & Ritual/rite & Endurance/leadership \\\bottomrule
\end{tabular}
\caption{Pillars}
\label{tab:pillars}
\end{table}

\paragraph{The Pillar of Severity}


The left pillar is feminine because pure intellect is the reflector of the Divine Ideas in Sophia. The pure intellect is based on compulsion, since its conclusions impinge on us by the force of their logical arguments. (It seems, however, that only mathematicians seem to ever actually be swayed by arguments.) In the practical life, its goal is Justice. The fundamental aspect of Justice is that a man gets what he deserves, or earns. Again, Justice is a form of compulsion. Finally, the willing function submits to the theories of the intellect and the notion of justice. It carries on the Tradition, especially the rites, rubrics, and so on.

The pure intellect needs to be invigorated by Wisdom, or else there is deterioration. No longer influenced by real ideas, it instead becomes enchanted by its own factitious and arbitrary thoughts. It follows the consequences of its formal logic with ineluctable certainty. Unable to intuit higher realities, the sense of justice is limited to material well-being. In \emph{The Two Sources of Morality and Religion}, \textbf{Henri Bergson} explains:

\begin{quotex}
If he were a slave of instinct, like the ant and the bee, he would remain intent on the purely external object to be attained; he would have automatically, somnambulistically, worked for the species. Endowed with intelligence, roused to thought, he will turn to himself and think only of leading a pleasant life. ... The truth is that intelligence would counsel egoism first. The intelligent being will rush in that direction if there is nothing to stop him.
\end{quotex}


Thus, in man's natural state, there is the choice between unconscious collectivism of blind instinct and materialistic individualism supported by an intellectuality ignorant of higher realities. The intellect necessarily operates on the known, or what is repeatable.

\begin{quotex}
Of the discontinuous alone does the intellect form a clear idea. ... Of immobility alone does the intellect form a clear idea. The intellect lets what is \emph{new} in each moment of history escape. It does not admit the unforeseeable. It rejects all creation. ... The intellect is characterised by a natural inability to comprehend life.
\flright{\textsc{Henri Bergson}}
\end{quotex}


Can anyone seriously doubt this? What do the finest scientific minds tell us? That there is no ``life'', just a complex chemical reaction. And there are no thoughts, but only electro-biochemical reactions in the brain. That ``artificial'' intelligence will be able to perform any task, better than humans. And ultimately the Turing test\footnote{\url{https://gornahoor.net/?p=8390}} will prove the impossibility of distinguishing an android from a human being. You can swim in that stagnant pond of pseudo-intellectuality if you like.

\paragraph{The Pillar of Mercy}


Bergson opposes instinctual knowledge to intellectual knowledge:

\begin{quotex}
Instinct is moulded on the very form of life. While intelligence treats everything mechanically, instinct proceeds, so to speak, organically. If the consciousness that slumbers in it should awake, if it were wound up into knowledge instead of being wound off into action, if we could ask and it could reply, it would give to us the most intimate secrets of life
\end{quotex}


Hence, the life of pure intellect becomes sterile without taking instinct into account. There is a natural law which is just known to be, either because it is innate or infused by a transcendent source. The intellect attempts to attack it, in order to replace it with an abstraction. However, unlike the ants and bees, instinct rises to something higher in man, what Bergson calls ``intuition'', or we are calling ``wisdom'':

\begin{quotex}
But it is to the very inwardness of life that intuition leads us ---by intuition I mean instinct that has become disinterested, self-conscious, capable of reflecting upon its object and of enlarging it indefinitely.
\end{quotex}


For example, the automated behaviour of the social insects rises to something higher in the human kingdom. The collective is served consciously and organically through custom, convention, continuity, and so on. That means that forced collectives, based on abstract logic and political power are rejected, even if the pure intellect believes them to be more efficient or sensible. Instead, the real collective is built up, starting with the family, then clan, tribe, and nation. These instinctual or intuited arrangements are considered to be divinely inspired in the sense that they transcend logical thought which can neither justify nor even comprehend them. Nevertheless, they are more certain since they are based on life itself, while intellectual movements come and go, they vacillate and can never overcome doubt.

Hence, instinct/intuition generates a morality higher than that of the formal intellect. Loyalty, fidelity, and piety become important, since disloyalty and impiety threaten the social arrangement. Loyalty is not to some abstract theory but to concrete entities: one's father and ancestors, one's family, nation, and religion. The virtues of the right pillar must be freely chosen. Hence, a man's loyalty is freely given, and obedience is voluntarily assumed. Now in the intellectual life a man can ``convert'' from one system to another, based on the ``evidence''. Such a man is inherently untrustworthy. On the other hand, to transfer loyalty from one group to another is treasonous.

Therefore, trustworthiness is proved by loyalty and obedience. A superior does not and cannot convince his subjects with ``evidence'' for his every decision; rather he must rely on the power of command. Nowadays, inferiors take notes to entrap their superiors, and loyalty is considered to be obstruction of justice.

At the feeling level, wisdom manifests as magnanimity, mercy, loving-kindness, and so forth. Unlike the corresponding ideal of justice, magnanimity is freely given as it is undeserved or unearned by the recipient. There is often a greater good to be achieved, but the sage must be disinterested in its application.

This ideal is poorly understood, and is usually distorted with a very modern notion of love. Unlike magnanimity, which is disinterested, active, and intelligent, this pseudo-love is subjective, passive, and instrumental. Magnanimity forgives despite the acts or character of the recipient, but its opposite holds that love needs to tolerate and even approve of the recipient. This pseudo-love is not disinterested, since it leads to strong emotions. It becomes instrumental, in that it claims that ``love'' will solve some social ill, etc. This sort of love is indiscriminate, whereas Tradition teaches that there is a hierarchy of love: self, wife, children, parents, siblings, friends, neighbours, fellow-countrymen, and then all others.

At the level of the willing function, this becomes endurance, strength, steadfastness, long-suffering until victory is achieved. Whenever the mind says to stop, surrender, or give up, the will nevertheless continues on.

\paragraph{Organic Logic}


In order to serve the common good, Individualism needs to move beyond purely formal logic into organic logic and moral logic\footnote{\url{https://www.meditationsonthetarot.com/moral-logic}}. While formal logic is concerned with the good of the individual, organic logic is concerned with the well-being of the collective. Social policies that sound rational and beneficial within formal logic, may collectively lead to adverse effects. Formal logic is quantitative, and so tends to regard human beings as fungible, in that qualities that would seem to distinguish human beings from each other are regarded as arbitrary and insignificant. Hence, policies are promoted on that basis, which, from a short term perspective, seem to aid the individual, yet eventually have deleterious effects on the family, nation, or religion. For example, in most ``advanced'' countries, the emphasis on the self-actualization of the individual leads to birth rates which are below replacement level with uncertain long term consequences. There are many other examples which could have been included. Obviously, this is a vast topic which cannot be addressed here.


\flrightit{Posted on 2017-06-11 by Cologero}

\begin{center}* * *\end{center}

\begin{footnotesize}
\begin{sffamily}
\texttt{James on 2017-06-12 at 08:44 said:}

Thanks so much for another great post.

\hfill

\texttt{TeenyTinyTarot on 2018-06-11 at 09:59 said:}

You write:

``This `magical enchantment' of following every ``impulse of the heart'' or any idea that pops into one's mind as guides to life . . . prevents the moral imagination from cognizing higher realities. This sequence that begins with sense life, which then excites the passions, and finally directing the mind, is a deformation of the normal sequence: viz., thought awakens feeling and directs the will.'' What do you make of the contrast between the ``normal sequence'' and that which characterizes ``The Hanged Man'' (aka the ``spiritual man'') in Letter XII:

``Now, the normal relationship between thought, feeling and the will for a civilised and educated man is such that his thought awakens feeling and directs the will. Thought plays a stimulating role, by means of imagination, towards feeling, and an educative role, by means of imagination and feeling, towards the will. Having to act. one thinks, one imagines, one feels, and ---lastly---one desires and acts. This is not so for the ``spiritual man''. For him it is the will which plays the stimulating and educative role towards feeling and thought. He acts first, then he desires, then he feels the worth of his action, and lastly he understands'' (The Hanged Man 316-317).

\end{sffamily}
\end{footnotesize}

